\section{Related Work}
\label{sec:related}

This section surveys recent research across four dimensions: traditional QoS routing, DRL-based routing, GNN-based network modeling, and graph network conditioning mechanisms.

\subsection{Traditional QoS Routing and Traffic Engineering}

Classical mechanisms such as OSPF, ECMP, and CSPF perform hop-by-hop forwarding based on static link weights and lack differentiated service capability under mixed traffic scenarios. Wang and Crowcroft established the constraint optimization framework for QoS routing~\cite{Wang1996_QoSRouting}; Fortz and Thorup revealed the limitations of weight-optimization-based traffic engineering under dynamic demand~\cite{Fortz2000_TrafficEngineering}. Brundiers et al.\ studied the benefits of loops in Segment Routing for traffic engineering, demonstrating the potential of flexible path control in dynamic load balancing~\cite{Brundiers2021_Loops_SR_TE}, and further proposed a midpoint optimization approach for finer-grained traffic engineering control~\cite{Brundiers2022_Midpoint_SR_INFOCOM}\textcolor{red}{这篇文献编译出来还是不对，check一下，找了一下原因，因为有另外一篇文献跟这个作者相同，所以被省略了，找下解决方案}. The common shortcoming of these methods lies in their myopia: each decision is independently optimized and cannot reserve resources for heterogeneous flows from a global network perspective or a long-term sequential view. When mouse flows and elephant flows share links, static strategies often fail to simultaneously satisfy strict latency guarantees and high throughput.

\subsection{DRL-Based Network Routing}

To overcome the limitations of fixed rules, researchers have introduced DRL into routing decisions. Liu et al.\ proposed DRL-OR, using multi-objective DRL to handle heterogeneous service demands~\cite{Liu2021_DRLOr_INFOCOM}. Casas-Velasco et al.\ proposed DRSIR, achieving intelligent routing decisions in SDN via deep Q-learning~\cite{CasasVelasco2021_DRSIR_TNSM}. Chen et al.\ proposed an adaptive multipath routing optimization framework based on multi-agent meta-reinforcement learning~\cite{Chen2022_MultiagentMeta_TNNLS}. Suzuki et al.\ proposed a multi-agent deep reinforcement learning method for jointly optimizing computation offloading and routing decisions in multi-cloud edge networks, validating the adaptive capability of multi-agent collaboration in dynamic network environments~\cite{Suzuki2023_MADRL_Routing_TNSM}. Sanchez et al.\ proposed DQS, achieving routing optimization in SDN via QoS-driven DRL~\cite{Sanchez2024_DQS_JPDC}. However, the above methods share two common limitations. First, their policy networks mostly encode network state via fully connected layers or simple embeddings, lacking explicit modeling of topological structure and thus suffering from limited generalization. Second, they typically optimize a single global objective for all flows (e.g., minimizing maximum link utilization or average latency) without embedding per-flow QoS intent as a conditioning signal into the policy network, and therefore cannot achieve differentiated routing under mixed traffic.

\subsection{GNN-Based Network Modeling and Routing}

GNNs provide a natural tool for encoding network topology, with existing research proceeding along two directions.

The first is offline performance modeling. Rusek et al.\ proposed RouteNet~\cite{Rusek2020_JSAC_RouteNet} and Ferriol-Galm\'{e}s et al.\ proposed RouteNet-Fermi~\cite{FerriolGalmes_2023_RouteNetFermi}, using message-passing neural networks to predict link delay and packet loss. Huang et al.\ proposed xNet, achieving flow-level QoS prediction by learning temporal state transitions~\cite{Huang2024_ToN_xNetModeling}. These methods focus on offline modeling under aggregated traffic, learning ``one graph, one global embedding,'' and cannot adapt to online per-flow decisions.

The second is graph-aware policies for online decision-making. Almasan et al.\ validated the improvement in routing quality from the joint application of GNNs and DRL~\cite{Almasan2022_RLmeetsGNN}; Bern\'{a}rdez et al.\ proposed MAGNNETO, a multi-agent GNN system for distributed traffic engineering~\cite{Bernardez2023_MAGNNETO_TCCN}; Yang et al.\ and Dong et al.\ respectively explored graph-aware joint optimization of routing and scheduling in time-sensitive networking and network slicing scenarios~\cite{Yang2022_GCN_RL_IoTJ,Dong2022_GCN_Slicing_Routing_TCCN}; Bhuvanasi et al.\ achieved resilient routing using GNNs and multi-agent DRL~\cite{Bhuvanasi2023_ResilientRouting_TNSM}. Although these methods introduce GNNs into online decision-making, they still produce a unified topological representation for all flows. When mouse flows and elephant flows arrive sequentially, the system sees the same embedding of the same graph, unable to perform differentiated perception based on per-flow QoS characteristics, let alone make different routing decisions accordingly. Collectively, these methods constitute the paradigm that couples GNN topology encoding with DRL path selection. The absence of flow conditioning makes this paradigm a natural but undifferentiated baseline.

\subsection{Conditioning Mechanisms for GNNs}

To enable network representations to dynamically adjust according to external conditions, researchers have proposed various conditioning mechanisms. These mechanisms can be classified by the level at which the conditioning signal operates into three categories:
(1)~\textbf{Node-level conditioning}, such as GNN-FiLM proposed by Brockschmidt et al.~\cite{Brockschmidt2020_GNNFiLM}, which generates affine modulation parameters from target node features to directly reshape the message-passing process; (2)~\textbf{Edge-level dynamic weights}, such as GATv2 by Brody et al.~\cite{Brody2022_GATv2}, which computes adaptive attention weights via dynamic interaction of node features, achieving conditioned adjustment of relational strength; (3)~\textbf{Type-level conditioning}, such as the Heterogeneous Graph Transformer by Hu et al.~\cite{Hu2020_HGT}, which uses node type as a semantic condition to differentiate feature transformations across different relation patterns.

However, the conditioning signals in the above mechanisms all originate from intra-graph structure (target nodes, edge features, node categories, etc.), rather than from per-flow external QoS intent. In the online routing scenario, what distinguishes embeddings is the QoS characteristics of the arriving flow: the same topology, when facing a mouse flow versus an elephant flow, should produce markedly different node representations and routing decisions. No prior work has introduced FiLM conditioning into the per-flow online routing scenario, and differentiated topology perception and policy learning remain open problems.

In summary, the existing literature exhibits clear gaps: (1)~traditional methods are myopic and cannot handle mixed QoS; (2)~DRL methods lack explicit topological modeling; (3)~GNN methods assume homogeneous traffic and suffer from the ``same graph, same embedding'' representational homogeneity problem; (4)~conditioned GNNs have not been applied to per-flow online routing. QoS-FiLM addresses these limitations.
